\documentclass[10pt,a4paper]{amsart}
\usepackage[margin=19mm]{geometry}
\usepackage{amsmath,amssymb,mathtools}
\usepackage[scr=rsfs,cal=euler]{mathalfa}
\usepackage{xfrac}
\usepackage[unicode,hidelinks,bookmarks=false]{hyperref}
\newcommand{\ie}{i.e.\ }
\newcommand{\cf}{cf.\ }
\newcommand{\resp}{respectively\ }
\newcommand{\Subsection}{\subsection}
\providecommand{\nobreakdash}{\nobreak\hbox{-}}
\setlength{\emergencystretch}{2em}
\multlinegap0pt
\usepackage{bm}
\usepackage{bbm}
\usepackage{dsfont}
\usepackage{xspace}
\usepackage{cancel}
\usepackage{mathtools}
\usepackage{amsthm}
\makeatletter
\@ifundefined{theorem}{\newtheorem{theorem}{Theorem}}{}
\@ifundefined{corollary}{\newtheorem{corollary}[theorem]{Corollary}}{}
\makeatother
\newcommand{\BB}{\mathbb{B}}
\newcommand{\FF}{\mathbb{F}}
\newcommand{\cI}{\mathcal{I}}
\newcommand{\cS}{\mathscr{S}}
\newcommand{\cT}{\mathcal{T}}
\newcommand{\cV}{\mathcal{V}}
\newcommand{\cW}{\mathcal{W}}
\newcommand{\tr}{{\rm Tr}}
\title[Trace Repair Never Loses to Classical Repair: Exact and Expl]{Trace Repair Never Loses to Classical Repair: Exact and Explicit Helper Nodes Selection}
\author{Wilton Kim, Stanislav Kruglik, Han Mao Kiah}
\date{Source: arXiv:2509.06492v2 (CC BY 4.0)}
\begin{document}
\maketitle

\noindent\emph{Source and licence.} Trace Repair Never Loses to Classical Repair: Exact and Explicit Helper Nodes Selection. Version arXiv:2509.06492v2 (CC BY 4.0), distributed under the Creative Commons Attribution 4.0 International licence, as recorded in the arXiv licence metadata for this exact version and shown on the abstract page https://arxiv.org/abs/2509.06492v2. Full rights notice: licenses/arxiv-2509.06492-CC-BY-4.0.txt.\par\smallskip

\noindent\textit{Editorial scope.} Counted unit: the Corollary whose source label is \texttt{cor:nodes\_to\_download} in the article (source0.tex lines 884--887, source label \texttt{cor:nodes\_to\_download}), delivered together with the article's own complete original proof of it (source0.tex lines 888--921, one \texttt{proof} environment of 34 lines). No mathematics is written, repaired or re-derived: statement and proof are the article's own text line for line. Required context retained uncounted: none. Every definition, display and label of those lines is retained unchanged. Omitted from this package and not redistributed: 1--883, 922--1072. Nothing inside a retained excerpt is omitted or altered: every retained body is delivered exactly as the article's lines are written, so no omission receipt of any kind accompanies this package. Citations: the retained excerpts carry no citation command at all, so no bibliography entry of the article is transcribed and no reference is redistributed. Reference adaptation: none was needed and none was made; every reference inside the delivered unit resolves inside it, so no unresolved reference or citation is produced. Printed slips kept literal: no printed word, symbol, hypothesis or number of the counted statement or of its proof was altered. Layout and class: the article's own class is \texttt{IEEEtran} with packages including amsmath, amsfonts, algorithmic, algorithm, array, subcaption, textcomp, stfloats, url, verbatim, graphicx, cite, comment, setspace; the wrapper is the lane's pinned \texttt{amsart} editorial wrapper at 10pt on A4, which restates the theorem-like environments the retained text needs and adds the article's own macro definitions that the retained text needs (\texttt{\textbackslash BB}, \texttt{\textbackslash FF}, \texttt{\textbackslash cI}, \texttt{\textbackslash cS}, \texttt{\textbackslash cT}, \texttt{\textbackslash cV}, \texttt{\textbackslash cW}, \texttt{\textbackslash tr}), with extra wrapper packages (bm, bbm, dsfont, xspace, cancel, mathtools). The delivered numbering is the unit's own. Assets: no retained span contains a figure environment, a graphics-inclusion command, a PDF-page inclusion, a table or a TikZ picture; no image file is copied into this package, the package declares no asset dependency and no image file is redistributed with it. Licence: version arXiv:2509.06492v2 of this article is distributed under the Creative Commons Attribution 4.0 International licence, as recorded in the arXiv licence metadata for this exact version and shown on the abstract page https://arxiv.org/abs/2509.06492v2. A full rights notice is delivered at licenses/arxiv-2509.06492-CC-BY-4.0.txt. Characters: every retained excerpt is pure ASCII, so the delivered engine, XeLaTeX with its default Latin Modern text font, reports no missing character.
\begin{corollary}\label{cor:nodes_to_download}
    With the output $(\Xi_{k,|\cS|}^*,|\cS|,R_{\cS})$ from Algorithm~O, consider basis element $\cT_{k,\cS}^*$ for $\cV_{k,\cS}$. Let $d = \dim(\cV_{k,\cS}) = \dim(\cW_{k,\cS})$. Fix $r$ and set $\cI = \{\omega^r,\ldots,\omega^{r+d-1}\}$ and $\cS$ to be any subset of $\FF^*\setminus\cI$ of the given size. Then, by downloading $\tr(g(\alpha)f(\alpha)/\alpha)$ for all $\alpha\in\FF^*\setminus(\cI\cup\cS)$, it is possible to recover $\tr(g(\alpha)f(\alpha)/\alpha)$ for all $\alpha\in\cI$. 
    Hence, we repair $f(0)$ with bandwidth $(n-1-R_\cS)\log|\BB|$ bits.
\end{corollary}

\begin{proof}
    To prove this, we refer back to the proof of Theorem~1. Note that, to perform the repair, we only need to ensure that the matrix $\boldsymbol{M}_\cI$ is invertible with some choice of $\cI$. Indeed, with $\cI = \{\omega^r,\ldots,\omega^{r+d-1}\}$, we have
    \begin{align*}
    \boldsymbol{M}_\cI&= \boldsymbol{VE} = \begin{bmatrix}
        \boldsymbol{V}_1 &\boldsymbol{0} & \cdots & \boldsymbol{0}\\
        \boldsymbol{0} & \boldsymbol{V}_2 & \cdots & \boldsymbol{0}\\
        \vdots & \vdots & \ddots & \vdots\\
        \boldsymbol{0} & \boldsymbol{0} & \cdots & \boldsymbol{V}_{|\Xi_k^*|}
    \end{bmatrix}\boldsymbol{E},
    \end{align*}
    where
    \begin{align*}
    \boldsymbol{E} = \begin{bmatrix}
        \left(\omega^{a^{(1)}}\right)^r & \cdots & \left(\omega^{a^{(1)}}\right)^{r+d-1}\\
        \vdots & \ddots & \vdots \\
        \left(\omega^{a^{(1)}q^{s_1-1}}\right)^r & \cdots & \left(\omega^{a^{(1)}q^{s_1-1}}\right)^{r+d-1}\\
        \vdots & \ddots & \vdots\\
        \left(\omega^{a_{1}^{(|\Xi_k^*|)}}\right)^r & \cdots & \left(\omega^{a_{1}^{(|\Xi_k^*|)}}\right)^{r+d-1}\\
        \vdots & \ddots & \vdots \\
        \left(\omega^{a^{(|\Xi_k^*|)}q^{s_{|\Xi_k^*|}-1}}\right)^r & \cdots & \left(\omega^{a^{(|\Xi_k^*|)}q^{s_{|\Xi_k^*|}-1}}\right)^{r+d-1}
    \end{bmatrix},
    \end{align*}
    and
    \begin{align*}
        \boldsymbol{V}_i = \begin{bmatrix}
            1 & \omega^{a^{(i)}} & \left(\omega^{a^{(i)}}\right)^2 & \cdots & \left(\omega^{a^{(i)}}\right)^{s_i-1}\\
            1 & \omega^{a^{(i)}q} & \left(\omega^{a^{(i)}q}\right)^2 & \cdots & \left(\omega^{a^{(i)}q}\right)^{s_i-1}\\
            \vdots & \vdots & \vdots & \ddots & \vdots\\
            1 & \omega^{a^{(i)}q^{s_i-1}} & \left(\omega^{a^{(i)}q^{s_i-1}}\right)^2 & \cdots & \left(\omega^{a^{(i)}q^{s_i-1}}\right)^{s_i-1}\\
        \end{bmatrix}. 
    \end{align*}
    Clearly, both $\boldsymbol{V}$ and $\boldsymbol{E}$ are invertible. This is because $\boldsymbol{V}$ is a block-diagonal matrix of Vandermonde matrices and $\boldsymbol{E}$ is a Vandermonde matrix. Then, the result follows. 
    %\hm{$\mathbfsl{V}$ and $\mathbfsl{E}$ should be Vandermonde matrices.}
\end{proof}

\end{document}
